\chapter{Pendahuluan}
\label{chap:pendahuluan}

\section{Latar Belakang}
\label{sec:latar-belakang}
Indonesia merupakan produsen minyak kelapa sawit mentah (CPO) terbesar di dunia
dengan kontribusi lebih dari 54\% ekspor global dan produksi sekitar 47 juta ton pada 2023
(Benedict dkk., 2024). Sektor ini menyumbang 3,5\% terhadap PDB nasional dan 13,5\% dari
ekspor non-migas (Kementerian Pertanian, 2022), serta menjadi penopang utama
perekonomian Indonesia. Konsumsi domestik juga meningkat dari 18,42 juta ton pada 2021
menjadi 23,13 juta ton pada 2023 (GAPKI,2024). Hal ini menandakan bahwa kelapa sawit telah bertransformasi menjadi komoditas yang sangat krusial, tidak hanya sebagai sumber devisa utama negara, tetapi juga sebagai penyangga stabilitas pasar dalam negeri yang permintaannya terus bertumbuh.
Namun, peningkatan produksi lebih banyak disebabkan oleh ekspansi lahan, bukan
peningkatan produktivitas per hektar (Abdul dkk., 2022). Salah satu penyebabnya adalah
menurunnya populasi kumbang penyerbuk \emph{Elaeidobius kamerunicus}, yang sebelumnya mampu
meningkatkan \emph{fruit set} lebih dari 20\% sejak introduksinya pada 1980-an (Azlan dkk., 2023).
Penurunan populasi kumbang disebabkan oleh beberapa faktor seperti cuaca ekstrem dan
infeksi nematoda yang dilaporkan menginfeksi hingga 91\% kumbang di beberapa wilayah
Malaysia dan Indonesia (Nur Farah dkk., 2025). Hal ini memiliki dampak langsung terhadap
efisiensi penyerbukan dan hasil tandan buah segar.
Berbagai inovasi telah dikembangkan untuk meningkatkan efisiensi perkebunan, seperti
sistem deteksi pohon sawit berbasis citra dengan akurasi tinggi (Putra dkk., 2022; Lee dkk.,
2024; Haq dkk., 2024), namun hingga saat ini belum ditemukan publikasi ilmiah yang secara
eksplisit membahas deteksi masa polinasi dan penyerbukan otomatis kelapa sawit. Pendekatan
\emph{hatch and carry} yang bertujuan untuk meningkatkan populasi kumbang penyerbuk telah
terbukti mampu meningkatkan \emph{fruit set} kelapa sawit sekitar 15–21\%, namun penerapannya di
lapangan masih terbatas oleh faktor operasional dan kondisi lingkungan (Gintoron dkk., 2023).
Maka dari itu, diperlukan adanya metode solutif untuk membantu kinerja kumbang
penyerbuk untuk melakukan penyerbukan kelapa sawit. Pengembangan alat pendeteksi masa
polinasi beserta sistem penyerbuk yang dapat memonitor volume ketersediaan \emph{pollen} dapat
menjadi solusi terkait menurunnya produktivitas kelapa sawit. Alat akan menangkap gambar
tandan bunga secara \emph{real-time} untuk dianalisis, dan jika terdeteksi bunga siap diserbuki, sistem
mengaktifkan mekanisme penyerbukan. Dengan ini, diharapkan dapat meningkatkan efisiensi
penyerbukan dan produktivitas kelapa sawit secara berkelanjutan.


\section{Rumusan Masalah}
\label{sec:rumusan-masalah}

Berdasarkan hal yang telah dipaparkan di latar belakang, berikut merupakan rumusan
masalah yang akan diselesaikan pada penelitian ini:
\begin{enumerate}
    \item Bagaimana cara alat dapat mendeteksi bunga betina kelapa sawit yang siap penyerbukan?
    \item Bagaimana desain dari alat penyerbuk yang dapat membantu penyerbukan pohon kelapa sawit?
    \item Bagaimana cara mengintegrasikan mekanisme penyerbuk dan sistem monitoring \emph{pollen} dengan \emph{drone} sebagai \emph{platform} pembawa?
\end{enumerate}

\section{Batasan Masalah}
\label{sec:batasan-masalah}

Ruang lingkup/pembatasan masalah dalam upaya memfokuskan penelitian yang akan
dilakukan menjadi lebih terarah.
\begin{enumerate}
  \item Proses penyerbukan hanya dilakukan pada pohon kelapa sawit dengan fokus pada bunga betina pada fase tertentu (sebelum dan saat reseptif), sehingga tidak mencakup pemodelan lengkap seluruh siklus fenologi tanaman mulai dari penyerbukan hingga buah panen.
  \item Desain alat penyerbuk dibatasi pada prototipe mekanis/elektromekanis yang mampu melakukan pelepasan serbuk atau media penyerbukan secara terarah, tanpa mengoptimalkan aspek \emph{industrial-scale} (daya tahan jangka panjang, standar industri, sertifikasi).
  \item Sistem monitoring volume \emph{pollen} dibatasi pada deteksi perubahan level \emph{pollen} menggunakan sensor tertentu, tanpa pengembangan metode prediksi lanjutan.
  \item Pengujian dilakukan pada luasan kebun terbatas dan dalam kondisi lingkungan tertentu (cuaca, kerapatan tanaman, topografi relatif sederhana), sehingga hasil tidak langsung digeneralisasi ke seluruh kondisi kebun kelapa sawit di Indonesia.
\end{enumerate}

\section{Tujuan}
\label{sec:tujuan}

Tujuan penelitian ini adalah sebagai berikut:
\begin{enumerate}
  \item Merancang dan mengimplementasikan metode penyerbukan pohon kelapa sawit yang terinspirasi dari perilaku \emph{Elaeidobius kamerunicus}, dengan memanfaatkan sensor citra untuk mengidentifikasi bunga betina yang siap diserbuki.
  \item Merancang dan membuat desain alat penyerbuk yang dapat diintegrasikan pada \emph{drone}, yang mampu mengaplikasikan serbuk atau media penyerbukan ke bunga kelapa sawit secara terarah dan terkendali.
  \item Mengevaluasi performa mekanisme penyemprotan serta akurasi deteksi volume \emph{pollen} pada kondisi pengujian nyata.
\end{enumerate}

\section{Manfaat}
\label{sec:manfaat}

Berikut adalah manfaat-manfaat yang dapat diambil dari penelitian ini:
\begin{enumerate}
  \item Memberikan kontribusi terhadap pengembangan kajian pemanfaatan sistem cerdas dalam proses penyerbukan kelapa sawit, khususnya pendekatan peniruan perilaku \emph{Elaeidobius kamerunicus} untuk deteksi kesiapan penyerbukan.
  \item Menjadi referensi awal untuk penelitian lanjutan terkait integrasi visi komputer pada aplikasi pertanian presisi di komoditas kelapa sawit.
  \item Membantu menyediakan prototipe solusi yang berpotensi mengurangi ketergantungan penuh pada penyerbukan manual, sehingga dapat mengurangi risiko penurunan produktivitas ketika populasi kumbang penyerbuk alami menurun atau tenaga kerja terbatas.
  \item Membuka peluang peningkatan efisiensi proses penyerbukan melalui penentuan target bunga yang lebih tepat (hanya yang siap diserbuki), sehingga penggunaan waktu, tenaga, dan serbuk lebih optimal.
  \item Menjadi dasar pengembangan produk/teknologi lanjutan yang dapat diadaptasi oleh perusahaan perkebunan atau peneliti lain untuk menguji dan mengembangkan sistem penyerbukan berbantuan robot secara lebih serius di tahap berikutnya.
\end{enumerate}

